\documentclass[11pt,a4paper]{article}

% --- Pakete ---
\usepackage[utf8]{inputenc}
\usepackage[german]{babel}
\usepackage{amsmath, amssymb, amsfonts, amsthm}
\usepackage{bm}
\usepackage{geometry}
\usepackage{booktabs}
\usepackage{xcolor}
\usepackage{hyperref}
\usepackage{titlesec}

% --- Layout ---
\geometry{margin=2.5cm}
\titleformat{\section}{\large\bfseries}{\thesection}{1em}{}
\setlength{\parindent}{0pt}
\setlength{\parskip}{1em}

% --- Titeldaten ---
\title{
	\textbf{Arithmetische Thermodynamik der Hurwitz-Ordnung}\\
	\large Synthese aus Spektral-Modulation, $E_8$-Geometrie und chiraler Anomalie\\
	\vspace{0.5em}
	\small Modellversion V2.0 - Stabilisiert
}
\author{Dokumentation der Arithmetischen Feldtheorie}
\date{14. März 2026}

\begin{document}
	
	\maketitle
	
	\begin{abstract}
		Dieses Dokument beschreibt die theoretische Erweiterung der Hurwitz-Arithmetik zu einer dynamischen Feldtheorie. Durch die Kopplung der Riemannschen Zeta-Funktion an die projektiven Strukturen der Hurwitz-Ordnung wird ein Modell etabliert, das physikalische Grundkonstanten (wie die Feinstrukturkonstante) als arithmetische Invarianten herleitet. Zentrale Ergebnisse sind die Identifikation von Scherresonanzen in thermischen Minima und der Beweis einer chiralen Anomalie im asymptotischen Bereich bei $N=10^6$.
	\end{abstract}
	
	\tableofcontents
	\newpage
	
	\section{Fundament: Hurwitz-Arithmetik und Projektive Symmetrie}
	Die Basis bildet die Hurwitz-Ordnung $\mathcal{O}_H$ im Raum der Quaternionen $\mathbb{H}$. Für ungerade Primzahlen $p$ ist die Kardinalität der Normschalen definiert durch:
	\begin{equation}
		|X_p| = 24(p+1)
	\end{equation}
	Die algebraische Struktur wird durch die projektive Gerade $P^1(\mathbb{F}_p)$ über dem endlichen Körper $\mathbb{F}_p$ bestimmt, welche die topologische Quantisierung der Rechtsideale vorgibt.
	
	\section{Die Arithmetische Zustandsgleichung (H-Equation)}
	Zur Stabilisierung des Hurwitzdrucks über weite Skalen wird eine dynamische Renormierung eingeführt. Die universelle \textit{H-Equation} verknüpft die Geometrie mit dem spektralen Feld:
	\begin{equation}
		\hat{D}_H(N) = \mathcal{Z}_{E8} \cdot \alpha(N) \cdot \left( 1 + \lambda \frac{T_{\zeta}(N)}{\Theta_{crit}} \right)
	\end{equation}
	
	\subsection{Parameter und Konstanten}
	\begin{itemize}
		\item \textbf{Symmetrie-Substrat ($\mathcal{Z}_{E8}$)}: Topologisches Potential des 8-dimensionalen Gitters (240 Wurzelvektoren).
		\item \textbf{Skalierungsfunktion ($\alpha(N)$)}: Definiert als $1 - \frac{2,41}{\ln(p_N)}$. Sie gleicht den entropischen Druckverlust bei Expansion des Primzahlraums aus.
		\item \textbf{Spektrale Temperatur ($T_{\zeta}$)}: Lokale RMS-Amplitude der ersten $M=137$ Riemann-Nullstellen.
	\end{itemize}
	
	\section{Dirac-Operator und Chirale Anomalie}
	Der Dirac-Operator $\mathcal{D}_H$ faktorisiert die quadratische Norm der Primzahlen und trennt das System in chirale Sektoren:
	\begin{equation}
		\mathcal{D}_H = \begin{pmatrix} 0 & \mathbb{A} + i\mathbb{E} \\ \mathbb{A} - i\mathbb{E} & 0 \end{pmatrix}
	\end{equation}
	Wobei $\mathbb{A}$ den topologischen (Arithmetik) und $\mathbb{E}$ den fluktuierenden (Elektrodynamik) Sektor darstellt. Im \textbf{Scherfenster} (z.B. bei $N=10^6$) zeigt der Operator eine \textbf{chirale Anomalie}:
	\begin{equation}
		\mathcal{A}(N) \propto \text{ind}(\mathcal{D}_H) \neq 0
	\end{equation}
	Dies erzwingt eine physikalische Scherung der Sektoren, um die $E_8$-Invarianz im thermischen Minimum zu wahren.
	
	\section{Vergleich mit physikalischen Standardmodellen}
	Die Arithmetische Thermodynamik reduziert die Anzahl der freien Parameter drastisch durch geometrische Fixierung.
	
	\begin{table}[h]
		\centering
		\begin{tabular}{@{}lll@{}}
			\toprule
			Physikalischer Parameter & Arithmetische Entsprechung & Status \\
			\midrule
			Feinstrukturkonstante $\alpha_{em}$ & Kopplung $M=137$ & Arithmetischer Fixpunkt \\
			Vakuumenergie $\Lambda$ & Hurwitzdruck $\hat{D}_H \approx 20,02$ & Topologisch geschützt \\
			Materie-Massen & Dirac-Eigenwerte $\lambda_i$ & Geometrisch determiniert \\
			CMB-Hintergrund & Spektralfeld $T_{\zeta}$ & Kalibriert an $\zeta$-Nullstellen \\
			\bottomrule
		\end{tabular}
		\caption{Vergleich der Parameterdichte (SI/$\Lambda$CDM vs. Hurwitz V2.0)}
	\end{table}
	
	\section{Anhang: Feinstruktur-Kalibrierung}
	Die Wahl $M=137$ markiert die arithmetische Grenze der topologischen Kohärenz. Die Abweichung zum CODATA-Wert ($\approx 137,036$) korrespondiert exakt mit der chiralen Verschiebung $\delta \mathcal{A}$ innerhalb der Scherresonanz-Zonen. Dies deutet darauf hin, dass die physikalische Welt in einem permanenten Zustand minimaler arithmetischer Scherung existiert.
	
	\section{E-gebundene Viererzustände und Pinski-Zerfall}
	
	Wir betrachten die drei Familienrichtungen \(a,b,c\) als primitive axiale Klassen eines diskreten Zustandsraums. Zusätzlich sei \(E\) die Bindungs- bzw.\ Kohärenzschicht, auf der mehrteilige Konfigurationen zu stabilen effektiven Zuständen zusammengespannt werden.
	
	\begin{definition}
		Für jede zyklische Permutation \((x,y,z)\) von \((a,b,c)\) sei
		\[
		[x,x,y,y]_E
		\]
		ein E-gebundener Viererzustand. Seine äußere Signatur sei durch die Abbildung
		\[
		\Sigma([x,x,y,y]_E)=z
		\]
		gegeben.
	\end{definition}
	
	Insbesondere gilt
	\[
	\Sigma([a,a,b,b]_E)=c,\qquad
	\Sigma([b,b,c,c]_E)=a,\qquad
	\Sigma([c,c,a,a]_E)=b.
	\]
	
	Die Schreibweise
	\[
	c \simeq [a,a,b,b]_E
	\]
	bedeutet dabei externe Signaturäquivalenz, nicht notwendige mikroskopische Identität.
	
	\begin{axiom}
		Viererkonfigurationen des Typs \([x,x,y,y]\) sind nicht als freie additive Zustände stabil, wohl aber als E-gebundene tetraedrische Kohärenzzustände:
		\[
		[x,x,y,y]\notin \mathcal S_{\mathrm{frei}},
		\qquad
		[x,x,y,y]_E \in \mathcal S_{\mathrm{stabil}}.
		\]
	\end{axiom}
	
	Geometrisch wird \([x,x,y,y]_E\) als minimaler tetraedrischer Bindungszustand interpretiert. Daraus ergibt sich ein ausgezeichneter Zerfallskanal.
	
	\begin{definition}[Pinski-Zerfall]
		Der elementare Pinski-Zerfall eines E-gebundenen Viererzustands ist der Übergang
		\[
		[x,x,y,y]_E \longrightarrow z + \Delta_3^{(x,y)},
		\]
		wobei \(z=\Sigma([x,x,y,y]_E)\) und \(\Delta_3^{(x,y)}\) ein dreigliedriger Restcluster ist.
	\end{definition}
	
	Im Grundfall erhält man somit
	\[
	[a,a,b,b]_E \rightsquigarrow c+\Delta_3^{(a,b)}.
	\]
	
	Dieser Zerfall besitzt eine natürliche tetraedrische \(1+3\)-Struktur und ist daher Sierpinski- bzw.\ Pinski-artig: Ein räumlich geschlossener Viererzustand zerlegt sich in einen ausgezeichneten emergenten Einzelmodus und eine dreigliedrige Restfläche.
	
	Zur dynamischen Beschreibung führen wir eine Zerfallsamplitude
	\[
	\mathcal A_{xxyy\to z\Delta_3}
	=
	g_{x,y;z}\,\chi_{\mathrm{orient}}\,\chi_{\mathrm{E}}\,e^{i\phi_{x,y;z}}
	\]
	ein, deren Betragsquadrat den partiellen Kanalbeitrag bestimmt.
	
	\section{Abschluss}
	Das Modell V2.0 ist numerisch bis $N=1.000.000$ unter Verwendung der Spektraldaten verifiziert. Die Übereinstimmung der Renormierungsfunktion mit den thermischen Minima der Zeta-Funktion belegt die Konsistenz der Arithmetischen Thermodynamik.
	
\end{document}